\section{Experimental setup}\label{sec:setup}

This section states what was run and how it is judged, in five parts:
Section~\ref{sec:setup-corpora} describes the three corpora and the label regime each
one represents, Section~\ref{sec:setup-split} the partitioning and why it is built around
participants rather than recordings, Section~\ref{sec:setup-train} the training protocol,
Section~\ref{sec:setup-metrics} the metrics and the statistical procedure, and
Section~\ref{sec:setup-repro} what one run consists of end to end together with the size
of the grid.

\subsection{Corpora and label provenance}\label{sec:setup-corpora}

A cough label can be obtained in at least three ways, and the three are treated here as
separate conditions rather than as interchangeable sources of ground truth. Table
\ref{tab:corpora} sets out the three corpora, what the positive and negative classes
consist of in each, and how the label was arrived at. The distinction is the organising
idea of the study, so it is worth stating plainly what each regime does and does not
establish.

\begin{table}[htbp]
\caption{The three corpora, the way the cough label was arrived at in each, and what the
two classes consist of. Every count is taken from the partition manifest the ingest
scripts write, not from the figures quoted by the corpus authors. The priors are matched
by construction so that a change in performance across regimes cannot be attributed to a
change in class balance.}\label{tab:corpora}
\small
\begin{tabular}{@{}lp{20mm}p{26mm}p{30mm}rrc@{}}
\toprule
Corpus & Label from & Positive class & Negative class & Recs. & Prior & Spk.\ ids \
\midrule
Coswara & a prompt & the two cough tasks, heavy and shallow & breathing, sustained vowel
  and counting tasks & $13{,}314$ & $0.2228$ & yes \
\addlinespace
FSD50K eval & a person & the \textsf{Cough} label & breathing, sneeze, gasp, sigh,
  laughter, speech & $476$ & $0.2227$ & no \
\addlinespace
COUGHVID v3 & a detector & \texttt{cough\_detected} $\geq 0.9$ &
  \texttt{cough\_detected} $\leq 0.1$ & $8{,}398$ & $0.2239$ & no \
\bottomrule
\end{tabular}
\end{table}

\textbf{Prompt derived labels.} Coswara~\cite{coswara2020} records each participant
performing a sequence of prescribed vocal tasks. A recording is labelled cough because
the participant was asked to cough into the microphone, not because anyone listened to
it. The positive class is the two cough tasks, heavy and shallow, and the negative class
is the breathing, sustained vowel and counting tasks. Two properties make this the corpus
the main experiments run on. It carries a participant identifier, which makes a speaker
disjoint partition possible, and both classes come from the same participant, device and
recording session, so a difference between the classes cannot be explained by channel or
speaker.

\textbf{Human annotated labels.} FSD50K~\cite{fsd50k2022} consists of material annotated
by people against the AudioSet ontology~\cite{gemmeke2017}, and its evaluation split was
exhaustively verified by the corpus authors. Only that split is used, and it is never
trained on. Its value here is twofold: the label means a human judged a cough to be
present, and the ontology provides negatives that are deliberately confusable, since
breathing, sneezing, gasping, sighing and laughter are separate classes that can be
selected for rather than left to chance. A detector that only separates coughing from
silence is not worth reporting, and this is the set that tests whether more than that has
been achieved. Other collections of human annotated vocal sounds exist and carry cough
among their classes~\cite{vocalsound2022}; this split was preferred because its labels
were exhaustively verified by the corpus authors rather than sampled for verification, and
because the released audio of the alternative could not be retrieved from its distribution
point at the time of writing.

\textbf{Classifier derived labels.} COUGHVID~\cite{coughvid2021} ships a
\texttt{cough\_detected} score produced by the corpus authors' own detector, and provides
no human cough or non-cough judgement for the corpus as a whole. A model trained and
tested on labels thresholded from that score is measured on its agreement with that
detector, not on whether a cough is present. This is reported as a property of the
regime rather than worked around, because it is the regime a large part of the published
literature on this corpus operates in.

Two further facts about COUGHVID are recorded here because they constrain what can be
concluded from it. The recordings carry anonymous identifiers and no participant field,
so a speaker disjoint partition cannot be constructed and any contributor who submitted
more than once may appear on both sides of a split. And the selection used by the work
being extended cannot be reproduced from the corpus as published. That work reports
$15{,}619$ cough and $7{,}261$ non-cough recordings, $22{,}880$ retained of $34{,}434$.
The lower figure is reproducible: a threshold of $0.1271$ on \texttt{cough\_detected}
yields exactly $7{,}261$. The upper figure is not reproducible by any threshold on that
column, because six recordings carry the value $0.8961$ exactly, so the count at or above
a threshold near there steps from $15{,}617$ directly to $15{,}623$ and $15{,}619$ lies
strictly between the two. Something further was applied that the original does not
record. Rather than tune a near miss, this work uses round thresholds and states them:
a score of $0.9$ or above is taken as cough and $0.1$ or below as non-cough, with the
interval between them discarded, which retains $22{,}003$ of $34{,}434$ recordings.

Because the two corpora differ sharply in class prior, the COUGHVID arm is subsampled to
the prior of the Coswara arm, keeping every non-cough recording and sampling the cough
class down. Without this the two regimes would differ in task difficulty as well as in
label provenance, and the comparison would not isolate what it is meant to isolate.

\subsection{Partitioning}\label{sec:setup-split}

The Coswara partition is built over participants, not recordings. Participants are
assigned to the training, validation and test partitions in the proportions $70$, $15$ and
$15$ per cent, stratified by reported status and by the number of cough recordings
contributed, under a fixed seed. Every recording from a participant follows that
participant, and the builder asserts that the three participant sets are pairwise
disjoint before it writes the manifest, so the property is enforced rather than intended.
One manifest serves all four conditions, so the conditions are compared on identical
partitions.

This is not a refinement. A partition made over recordings allows one contributor to
appear in both training and test material, which makes part of the test set a memory
test. The scale of the effect has been measured before: a mel spectrogram convolutional
network reported $98.18$ per cent accuracy under a random split and $95.18$ per cent when
test coughs came from recordings unseen in training~\cite{zhou2021}, a gap of roughly
three points which that work reported without discussing. Evidence from respiratory audio
more broadly points the same way. Audio classifiers that appeared accurate on
crowdsourced material fell from $0.846$ to $0.619$ area under the curve once confounders
including self-reported symptoms were matched~\cite{coppock2024}, and a careful evaluation
of the same kind of data reached $0.71$ where a less careful one reached
$0.90$~\cite{han2022}. A recent cross-dataset study of cough models found that their
learned representations cluster by recording device and corpus rather than by the
condition being detected~\cite{zhang2026}, which is the mechanism these gaps would be
expected to arise from. Two of its authors are also authors of
COUGHVID~\cite{coughvid2021}.

For COUGHVID the partition is over recordings, for the reason given in
Section~\ref{sec:setup-corpora}, and that limitation attaches to every number reported
for that arm. FSD50K is used only for evaluation, so it is not partitioned at all.

\subsection{Training protocol}\label{sec:setup-train}

One protocol is applied identically to every run. Each network is trained with
Adam~\cite{kingma2015} at a learning rate of $10^{-4}$ for at most eight epochs, with
gradients clipped to unit norm. Training stops early when validation balanced accuracy
has failed to improve for three consecutive epochs. The weights retained are those from
the epoch with the best validation balanced accuracy, selected without any reference to
the test partition, and the test partition is read exactly once per run, after selection.
Computation is in single precision. Table~\ref{tab:training} gives the full settings.

Three of these choices are consequences of observed failures rather than preferences.
Gradient clipping is present because an unclipped run at the learning rate inherited from
the work being extended diverges to a non-finite loss on this hardware. Single precision
is used because mixed precision produced a non-finite loss from the first optimiser step:
at the gradient scaler's initial scale the gradients of the first convolution overflow
the half precision range, and clipping then divides by an infinite norm and writes
non-finite values into the gradients after the scaler's own overflow check has already
run. The network never leaves the majority class, and the resulting accuracy equals the
majority class frequency exactly. The trainer counts and reports any optimiser step it
has to abandon for a non-finite gradient norm, so this mode of failure cannot pass for a
successful run.

A third choice follows from a measurement rather than a preference, and it concerns the
batch size, which is set per architecture. The driver on the development hardware does
not refuse an allocation that exceeds the four gigabytes of device memory. It satisfies
the excess from host memory instead, and the traffic then crosses the system bus on every
access. Nothing fails, no warning is issued, and the only symptom is throughput.
MobileNetV2 at a batch of $64$ reserves $5450$ MiB and trains at $23$ images per second,
while the same network at $16$ reserves $2228$ MiB and trains at $148$, a factor of
$6.4$ for a quantity that appears in no diagnostic. DenseNet201 shows the same effect
more sharply, $1.6$ images per second at a batch of $32$ against $29.8$ at $8$. Inverted
residuals and dense concatenation are what make these two the binding cases, since both
keep high resolution feature maps alive across many layers, so the activation memory of
each is far larger than its operation count would suggest. The competing explanation,
that the depthwise separable convolutions were reaching a slow kernel, was tested by
repeating every measurement in the channels last memory layout and is rejected: that
layout leaves the peak essentially unchanged and runs an order of magnitude slower again,
and a remedy that raises throughput without lowering the peak was not addressing a memory
problem. The batch for each architecture is therefore set to leave roughly a gigabyte of
the card unused rather than to the largest value that happens to fit, because the
quantity being protected is a peak and the rest of the machine is in use.
Table~\ref{tab:throughput} gives the measurements the choice rests on.

\begin{table}[htbp]
\caption{Training protocol. One protocol is applied identically to every run. The batch
size is the only setting that varies with architecture, for the reason given in the
text.}\label{tab:training}
\small
\begin{tabular}{@{}lp{62mm}@{}}
\toprule
Setting & Value \
\midrule
Optimiser & Adam, learning rate $10^{-4}$ \
Gradient clipping & unit norm \
Epochs & at most $8$ \
Early stopping & no gain in validation balanced accuracy for $3$ epochs \
Weights retained & those of the best validation balanced accuracy epoch \
Test reads per run & one, after selection \
Precision & single \
Input & $224 \times 224$, normalised by ImageNet channel statistics \
Batch size & $64$ ResNet18, $16$ MobileNetV2, $128$ compact, $8$ DenseNet201 \
Seeds & five, and three for DenseNet201 \
\bottomrule
\end{tabular}
\end{table}

\begin{table}[htbp]
\caption{Training throughput and peak reserved device memory, measured per architecture
on the card used throughout. The card holds $4096$ MiB. A peak above that figure is not
an error on this platform: the excess is served from host memory, and the collapse in
throughput is the only indication. The channels last rows show that the layout does not
address the cause.}\label{tab:throughput}
\small
\begin{tabular}{@{}llrrr@{}}
\toprule
Architecture & Layout & Batch & Images/s & Peak MiB \
\midrule
ResNet18 & contiguous & $64$ & $208.1$ & $2898$ \
\addlinespace
MobileNetV2 & contiguous & $64$ & $23.2$ & $5450$ \
MobileNetV2 & contiguous & $32$ & $148.4$ & $3222$ \
MobileNetV2 & contiguous & $16$ & $147.6$ & $2228$ \
MobileNetV2 & channels last & $32$ & $13.2$ & $3254$ \
MobileNetV2 & channels last & $16$ & $14.8$ & $2238$ \
\addlinespace
Compact & contiguous & $128$ & $967.1$ & $3098$ \
Compact & contiguous & $64$ & $912.6$ & $1608$ \
\addlinespace
DenseNet201 & contiguous & $32$ & $1.6$ & $7372$ \
DenseNet201 & contiguous & $12$ & $33.6$ & $3070$ \
DenseNet201 & contiguous & $8$ & $29.8$ & $2284$ \
\bottomrule
\end{tabular}
\end{table}

\subsection{Metrics and statistical procedure}\label{sec:setup-metrics}

Accuracy alone is not adequate on these class priors, so balanced accuracy, specificity
and the Matthews correlation coefficient~\cite{matthews1975} are reported alongside it.
The correlation coefficient is included because it responds to all four cells of the
confusion matrix in proportion to both class sizes, which accuracy and the harmonic mean
of precision and recall do not~\cite{chicco2020}. Interval estimates use the Wilson score
interval~\cite{wilson1927} rather than the Wald interval, whose coverage is erratic at
proportions near one~\cite{brown2001}, which is the region every condition here occupies.

Differences between conditions are tested on paired decisions, not on accuracies. For a
pair of conditions evaluated on the same test recordings under the same seed, the
quantities of interest are the counts of recordings that one condition classifies
correctly and the other does not, and these are assessed with McNemar's
test~\cite{mcnemar1947}, exactly while the discordant total is small enough to enumerate
and by the continuity corrected approximation otherwise. Of five tests compared for this
purpose, McNemar's was the only one with an acceptable Type I error rate for a learner
that is trained once~\cite{dietterich1998}. Because a family of such comparisons is made,
the p values are adjusted by the step down procedure of Holm~\cite{holm1979}.

One further quantity is reported beside every effect, and it governs how the effects
should be read. Each condition is run under five seeds, and the spread of accuracy across
those seeds is reported. An effect smaller than roughly twice that spread cannot be
claimed from single runs, however large the test set is. This is stated because
differences well under one accuracy point are routinely reported in this area without any
estimate of run to run variation, and the present work is in a position to say what that
variation is.

\subsection{Reproducibility}\label{sec:setup-repro}

A single run consists of the following, end to end. The images for one condition are
loaded for all three partitions and held in device memory as unsigned bytes, converted to
normalised single precision per batch. One architecture is trained from one seed under
the protocol of Section~\ref{sec:setup-train}. The epoch with the best validation
balanced accuracy is restored, the test partition is evaluated once, and one row per test
recording is written out carrying the file identifier, the true label, the predicted
label and both class scores. Those rows, not any summary of them, are what the analysis
reads, which is what makes the paired tests of Section~\ref{sec:setup-metrics} possible
at all.

The grid over the prompt derived corpus is $66$ runs. ResNet18, MobileNetV2 and the
compact network are each run on all four conditions under five seeds, which is $60$ runs,
and DenseNet201 is run on the two mel conditions under three seeds, which is a further
six; its batch size constraint makes it several times more expensive per run than the
others. The four conditions are the two frequency axes crossed with denoising applied or
withheld.

Three secondary experiments are reported, each settling one question.
\begin{enumerate}
\item Whether the accuracy obtained on prompted recordings survives a change of corpus,
      annotator and negative class. Settled by evaluating the retained networks on the
      human annotated evaluation split without retraining, in
      Section~\ref{sec:res-external}.
\item Whether the label regime changes the answer about denoising. Settled by repeating
      the ResNet18 block on the classifier derived corpus at a matched class prior, in
      Section~\ref{sec:res-regime}.
\item Whether the expert annotated subset of the classifier derived corpus can be used to
      estimate its label error rate. Settled in the negative, in
      Section~\ref{sec:res-regime}, because that subset is not a random sample of the
      corpus.
\end{enumerate}

Every count in this section and in Section~\ref{sec:results} is derived from the saved
decision files by script rather than by arithmetic on the page, and the scripts that do
so are released with the code.
